% \documentclass{article}
% \usepackage{graphicx} % Required for inserting images

% \title{Rubric for Homework 2}

% \begin{document}

% \maketitle

% \section{$Q_1$ Game Theory}
% \begin{itemize}
%     \item 1
% \end{itemize}


% \end{document}

\documentclass[11pt, letterpaper]{article}

% PACKAGES
\usepackage[margin=1in]{geometry}
\usepackage{amsmath}
\usepackage{amssymb}
\usepackage{graphicx}
\usepackage[utf8]{inputenc}
\usepackage{hyperref}

% DOCUMENT LAYOUT
\setlength{\parindent}{0pt}
\setlength{\parskip}{1em}
\renewcommand{\labelitemi}{\textbullet}
\renewcommand{\labelitemii}{$\circ$}

\hypersetup{
    colorlinks=true,
    linkcolor=blue,
    filecolor=magenta,      
    urlcolor=cyan,
}

% HELPER COMMANDS
\newcommand{\question}[2]{
    \subsection*{Question #1: #2}
}
\newcommand{\points}[1]{
    \textbf{(#1 Points)}
}

\begin{document}

% --- TITLE ---
\begin{center}
    {\Huge \textbf{Midterm Exam Grading Rubric}} \\
    \vspace{0.2cm}
    {\large IEDA 1250: Optimization for Decision Making (Summer, 2026)} \\
    \vspace{0.5cm}

    \hrule
\end{center}

% --- QUESTION 1 ---


% --- QUESTION 1 ---
\question{1}{Linear Programming -- Graphical Method \points{10}}

\begin{itemize}
    \item Correctly identifies that the feasible region is a diamond bounded by
    $0 \leq x_1+2x_2 \leq 6$ and $-2 \leq 2x_1-x_2 \leq 4$. \points{2}

    \item Correctly finds the four vertices of the feasible region:
    \[
    \left(-\frac{4}{5},\frac{2}{5}\right),
    \left(\frac{8}{5},-\frac{4}{5}\right),
    \left(\frac{2}{5},\frac{14}{5}\right),
    \left(\frac{14}{5},\frac{8}{5}\right).
    \]
    \points{4}

    \item Correctly identifies the optimal solution and optimal objective value:
    \[
    x_1^*=\frac{14}{5},\qquad x_2^*=\frac{8}{5},\qquad z^*=\frac{18}{5}.
    \]
    \points{4}
\end{itemize}


% --- QUESTION 2 ---
\question{2}{Linear Programming -- Complementary Slackness \points{20}}

\begin{itemize}
    \item Correctly writes or uses the dual constraints corresponding to the primal LP,
    taking into account that $x_3,x_4,x_5$ are unrestricted in sign. \points{4}

    \item Correctly applies complementary slackness to the primal constraints and obtains
    \[
    2x_1+x_2+x_3+x_4=3.
    \]
    \points{4}

    \item Correctly applies complementary slackness to the dual constraints for $x_1$
    and $x_2$, and concludes that
    \[
    x_1=0,\qquad x_2=0.
    \]
    \points{5}

    \item Correctly substitutes $x_1=0$ and $x_2=0$ into the primal constraints and forms
    the reduced system:
    \[
    \begin{cases}
    x_3+x_4=3,\\
    2x_3+x_5=7,\\
    x_3-x_4+2x_5=7.
    \end{cases}
    \]
    \points{4}

    \item Correctly solves the reduced system and obtains the optimal primal solution:
    \[
    x_1^*=0,\quad x_2^*=0,\quad x_3^*=2,\quad x_4^*=1,\quad x_5^*=3.
    \]

    the optimal objective value:
    \[
    z^*=6.
    \]
    
    \points{3}

    % \item Correctly reports the optimal objective value:
    % \[
    % z^*=6.
    % \]
    % \points{1}
\end{itemize}


\question{4}{Mixed-Integer Linear Programming -- Production Planning \points{20}}

\begin{itemize}
    \item Correctly defines production, inventory, and setup decision variables,
    including binary setup variables $y_t$ for each month. \points{4}

    \item Correctly formulates the objective function with production costs, fixed
    setup costs, and holding costs for Months 1--5. \points{4}

    \item Correctly writes the monthly inventory balance constraints using zero initial
    inventory and all monthly demands. \points{5}

    \item Correctly imposes no ending inventory, $I_6=0$, and nonnegativity of production
    and inventory variables. \points{3}

    \item Correctly links production and setup decisions using capacity constraints
    such as $0\leq x_t\leq 80y_t$ and states $y_t\in\{0,1\}$. \points{4}
\end{itemize}


\question{5}{Game Theory -- Dominance and Pure Strategy Equilibrium \points{20}}

\begin{itemize}
    \item Correctly recognizes that Player A maximizes the payoff and Player B minimizes
    the payoff in the zero-sum game. \points{2}

    \item Correctly eliminates dominated rows for Player A:
    Rows $A3,A4,A5$ are eliminated. \points{5}

    \item Correctly eliminates dominated columns for Player B:
    Columns $B3,B4,B5$ are eliminated. \points{5}

    \item Correctly obtains the reduced game
    \[
    \begin{array}{c|cc}
     & B1 & B2\\
    \hline
    A1 & 4 & 1\\
    A2 & 2 & 3
    \end{array}
    \]
    \points{3}

    \item Correctly checks the saddle point condition:
    $\max\min=2$ and $\min\max=3$, so no pure strategy Nash equilibrium exists.
    \points{3}

    \item Correctly explains why each payoff in the reduced game is unstable by
    identifying a profitable unilateral deviation by Player A or Player B. \points{2}
\end{itemize}


\question{6}{Integer Programming -- Logical Constraints \points{30}}

\begin{itemize}
    \item Correctly defines binary project selection variables $x_i\in\{0,1\}$ and
    formulates the profit-maximization objective and budget constraint. \points{4}

    \item Correctly formulates constraints (a)--(h):
    \[
    x_2+x_5\leq 1,\quad
    \sum_{i=1}^{10}x_i=4,\quad
    x_1+x_3+x_4\geq 1,
    \]
    \[
    x_6\leq x_7,\quad
    x_8\leq x_3+x_4,\quad
    x_1=x_{10},
    \]
    \[
    x_9\geq x_3+x_4-1,\quad
    x_6+x_7+x_8\geq 2x_5.
    \]
    \points{14}

    \item Correctly models condition (i), including an auxiliary binary variable and
    valid linear constraints enforcing at least one of the two stated conditions.
    One valid formulation is
    \[
    x_1+x_2+x_3+x_4\geq 2z_i,
    \]
    \[
    x_5+x_6+x_7+x_8\leq 1+3z_i,\qquad z_i\in\{0,1\}.
    \]
    \points{5}

    \item Correctly models condition (j), including an auxiliary binary variable and
    valid linear constraints for the conditional disjunction. One valid formulation is
    \[
    x_2\geq x_{10}+z_j-1,\qquad
    x_9\geq x_{10}+z_j-1,
    \]
    \[
    x_3\leq 1-x_{10}+z_j,\qquad
    x_4\leq 1-x_{10}+z_j,\qquad
    z_j\in\{0,1\}.
    \]
    \points{5}

    \item Uses clear notation and states all required binary restrictions for auxiliary
    variables. \points{2}
\end{itemize}



\end{document}
